\section{Neural networks, from zero}\label{sec:nn}

\subsection{One neuron}

A \textbf{neuron} takes some numbers in, multiplies each by its own \textbf{weight}, adds everything
up, adds one more number called the \textbf{bias}, and passes the result through a simple rule
called \textbf{ReLU}: negative numbers become 0, positive numbers stay as they are.
\[
  \text{output} = \text{ReLU}\Big(\sum_k x_k\,w_k + b\Big),\qquad \text{ReLU}(v) = \max(v, 0).
\]
Each $x_k w_k$ is one multiplication, and adding it to the running sum is one
\textbf{multiply-accumulate}, or \textbf{MAC}. MACs are almost all the work a network does.

\begin{figure}[H]
\centering
\begin{tikzpicture}[>=Stealth,node distance=1.0cm,
  inp/.style={circle,draw,minimum size=0.8cm,fill=gray!8},
  neu/.style={circle,draw,thick,minimum size=1.3cm,fill=roba!15}]
  \node[inp] (x1) at (0,1.4) {2};
  \node[inp] (x2) at (0,0) {0};
  \node[inp] (x3) at (0,-1.4) {3};
  \node[neu] (n) at (4,0) {$\Sigma$, ReLU};
  \draw[->] (x1) -- node[above,sloped,font=\small] {$w=1$} (n);
  \draw[->] (x2) -- node[above,font=\small] {$w=-2$} (n);
  \draw[->] (x3) -- node[below,sloped,font=\small] {$w=4$} (n);
  \node[above=0.1cm of n,font=\small] {bias $b=1$};
  \draw[->] (n) -- ++(2.2,0) node[right] {15};
  \node[font=\small,anchor=west] at (7.4,0) {$2\cdot1 + 0\cdot(-2) + 3\cdot 4 + 1 = 15$};
\end{tikzpicture}
\caption{One neuron with three inputs. It does three multiplications. If the weights were
$(-3,1,1)$ and the bias 0, the sum would be $-6+0+3=-3$, and ReLU would output 0.}
\end{figure}

\begin{kid}
A neuron is like a judge giving points. Each clue (input) has an importance (weight). The judge adds
up \emph{clue $\times$ importance} for all clues. If the total is negative, the judge says ``nothing'' (0).
\end{kid}

\subsection{Layers, scores and the answer}

A \textbf{layer} is many neurons that all look at the same inputs. The outputs of one layer are the
inputs of the next. The last layer has one neuron per possible answer (10 here: the digits 0--9, or 10
kinds of clothes). Its outputs are called \textbf{scores} or \textbf{logits}, and the network's answer is
the class with the \emph{largest} score (``argmax'').

\begin{example}[: from scores to an answer, and the loss]
Suppose the scores for three classes are $(2,1,0)$. The answer is class 0. During training, scores are
turned into probabilities with \emph{softmax}: $e^2=7.39$, $e^1=2.72$, $e^0=1$, total $11.1$, so the
probabilities are $0.665, 0.245, 0.090$. If class 0 is correct, the \textbf{loss} is
$-\ln 0.665 = 0.41$. If the network were sure (probability 1), the loss would be 0. The loss measures how
``surprised'' the network is by the right answer. Training changes the weights to make the loss small.
\end{example}

\textbf{Training} means: show the network many pictures with the right answers, compute the loss, and
nudge every weight a tiny bit in the direction that lowers the loss (the direction is found by
\emph{backpropagation}; the nudging rule we use is called \emph{Adam}). One pass over all training
pictures is an \textbf{epoch}. \textbf{Test accuracy} is the share of \emph{new} pictures (never seen
in training) that the network gets right.

\subsection{The two datasets}

\begin{itemize}
  \item \textbf{MNIST}: 70{,}000 hand-written digits (0--9), each $28\times 28$ grey pixels; 60{,}000 for
        training and 10{,}000 for testing. Easy: good networks get about 99\%.
  \item \textbf{Fashion-MNIST}: the same size and format, but the 10 classes are T-shirt, trouser,
        pullover, dress, coat, sandal, shirt, sneaker, bag and ankle boot. Harder (about 90\%), because
        a shirt and a pullover can look very alike.
\end{itemize}

Each pixel is a number from 0 (black) to 255 (white). Before the network sees it, it is divided by 255
and then \textbf{normalised}: $x = (p/255 - \text{mean})/\text{std}$, with mean $0.1307$ and std
$0.3081$ for MNIST (mean $0.2860$, std $0.3530$ for Fashion-MNIST). Normalising keeps the numbers centred
around 0, which makes training easier.

\begin{example}[: normalising a pixel]
For MNIST, a black pixel ($p=0$) becomes $(0-0.1307)/0.3081 = -0.424$ and a white pixel ($p=255$)
becomes $(1-0.1307)/0.3081 = 2.82$. Note that black is \textbf{not zero} after normalising. This will
matter a lot in Section~\ref{sec:zeroskip}.
\end{example}

\subsection{Network 1: the MLP}

An \textbf{MLP} (multi-layer perceptron) is the simplest network: layers in which every neuron sees
every input (``fully connected'', FC). Ours flattens the $28\times28$ picture into 784 numbers and uses
three layers:
\[
  784 \xrightarrow{\text{FC}} 256 \xrightarrow{\text{ReLU, FC}} 128 \xrightarrow{\text{ReLU, FC}} 10.
\]

\begin{example}[: how many multiplications does the MLP do?]
Every neuron multiplies each of its inputs once:
\[
  784\times256 + 256\times128 + 128\times10 = 200{,}704 + 32{,}768 + 1{,}280 = \mathbf{234{,}752}
\]
multiplications per picture. 85\% of them are in the first layer.
\end{example}

\subsection{Network 2: LeNet-5, with convolutions}

A \textbf{convolution} layer slides a small grid of weights (a \textbf{filter}, here $5\times5$) over
the picture. At every position, it multiplies the 25 pixels under the filter by the 25 weights and adds
them up. The same filter is used everywhere, so it finds the same pattern (an edge, a curve) wherever it
appears.

\begin{example}[: a $3\times3$ filter that finds vertical edges]
Filter $\begin{pmatrix}1&0&-1\\1&0&-1\\1&0&-1\end{pmatrix}$.
On a patch with a bright left part and a dark right part,
$\begin{pmatrix}9&9&0\\9&9&0\\9&9&0\end{pmatrix}$, the result is $9+9+9-0-0-0 = 27$: an edge.
On a flat patch $\begin{pmatrix}9&9&9\\9&9&9\\9&9&9\end{pmatrix}$, the result is
$27-27 = 0$: no edge. Each position costs 9 multiplications.
\end{example}

\textbf{Max pooling} ($2\times2$) shrinks a picture by keeping the biggest of every 4 neighbours, e.g.
$\begin{pmatrix}1&3\\2&0\end{pmatrix}\to 3$. It has no multiplications.

\begin{figure}[H]
\centering
\begin{tikzpicture}[>=Stealth,font=\small,
  blk/.style={draw,rounded corners=2pt,align=center,minimum height=1.0cm,inner sep=3pt}]
  \node[blk,fill=gray!10] (in) {picture\\$1\times28\times28$};
  \node[blk,fill=roba!15,right=0.35cm of in] (c1) {conv1\\6 filters $5\times5$\\$6\times28\times28$};
  \node[blk,fill=gray!5,right=0.35cm of c1] (p1) {pool\\$6\times14\times14$};
  \node[blk,fill=roba!15,right=0.35cm of p1] (c2) {conv2\\16 filters $6\times5\times5$\\$16\times10\times10$};
  \node[blk,fill=gray!5,right=0.35cm of c2] (p2) {pool\\$16\times5\times5$\\$=400$};
  \node[blk,fill=roba!15,below=0.9cm of p2] (f1) {FC\\$400\to120$};
  \node[blk,fill=roba!15,left=0.35cm of f1] (f2) {FC\\$120\to84$};
  \node[blk,fill=roba!15,left=0.35cm of f2] (f3) {FC\\$84\to10$};
  \node[blk,fill=gray!10,left=0.35cm of f3] (out) {10 scores};
  \draw[->] (in) -- (c1); \draw[->] (c1) -- (p1); \draw[->] (p1) -- (c2); \draw[->] (c2) -- (p2);
  \draw[->] (p2) -- (f1); \draw[->] (f1) -- (f2); \draw[->] (f2) -- (f3); \draw[->] (f3) -- (out);
\end{tikzpicture}
\caption{LeNet-5 as used here (ReLU after every layer except the last).}
\end{figure}

\begin{example}[: how many multiplications does LeNet-5 do?]
\begin{tabular}{@{}ll@{}}
conv1: & 6 filters $\times$ $28\times28$ positions $\times$ 25 weights $= 117{,}600$\\
conv2: & 16 filters $\times$ $10\times10$ positions $\times$ $(6\times25=150)$ weights $= 240{,}000$\\
FC layers: & $400\times120 + 120\times84 + 84\times10 = 48{,}000+10{,}080+840 = 58{,}920$\\
total: & $\mathbf{416{,}520}$ multiplications per picture.
\end{tabular}
\end{example}

\begin{inthecode}{roba/nn.py}
\code{class MLP} and \code{class LeNet5} define the two networks. Their layers are \code{ALinear} and
\code{AConv2d}: normal PyTorch layers whose multiplications can be switched to any of our eight
multipliers and any number format (``domain'').
\end{inthecode}

\subsection{Turning decimals into whole numbers: quantization}\label{sec:quant}

A trained network's weights are decimals like $0.2$ or $-0.0371$. Our small CPU has no hardware for
decimals (no floating-point unit), and the approximate multipliers work on whole numbers anyway. So we
\textbf{quantize}: represent every decimal as a whole number times a fixed \emph{scale}.

For 8 bits (\textbf{int8}), the whole numbers go from $-127$ to $127$:
\begin{itemize}
  \item Weights: for each output neuron, $s_w = \max|w| / 127$ and $w_{\text{int}} = \text{round}(w/s_w)$.
  \item Activations (layer inputs): $s_x = (\text{largest input seen on 5{,}000 training pictures})/127$
        and $x_{\text{int}} = \text{round}(x/s_x)$.
  \item Then $\sum x_k w_k \approx s_x s_w \sum x_{\text{int},k}\, w_{\text{int},k}$. The sum is
        made of whole-number products, which is what the approximate multipliers compute.
\end{itemize}

\begin{example}[: quantizing one product]
Say the largest weight of a neuron is $0.5$, so $s_w = 0.5/127 = 0.003937$. A weight $w=0.2$ becomes
$\text{round}(0.2/0.003937) = \text{round}(50.8) = 51$.
Say the largest input is $4.0$, so $s_x = 4/127 = 0.0315$. An input $x=1.3$ becomes
$\text{round}(41.3) = 41$.
The integer product is $41\times 51 = 2091$, and $2091\times s_x s_w = 2091\times 0.000124 = 0.2593$.
The true product is $1.3\times0.2 = 0.26$. Very close.
\end{example}

\begin{careful}
Simply cutting decimals to integers \emph{without} a scale (``trunc'') destroys the network:
$\text{trunc}(0.2) = 0$, and nearly every weight is between $-1$ and $1$, so nearly every weight
becomes 0. The network then always gives the same answer and gets about \textbf{10\%}, the same as
guessing (Figure~\ref{fig:domains}). We kept this case in the results to show why scaling is needed.
\end{careful}

We also tried \textbf{int16} ($-32767\dots32767$), which is more precise. For these networks it gives
the same accuracy as int8, so the Ibex part uses int8.

\subsubsection*{Integer-only rescaling (what the CPU actually does)}

After a layer adds up its integer products into \code{acc}, the next layer needs its input as int8
again, with its own scale $s_{\text{next}}$. On the GPU we just multiply by
$s_x s_w/s_{\text{next}}$. On Ibex there are no decimals, so we use \textbf{fixed point}: store the
decimal scale as a whole number $M$ that is $2^{24}$ times bigger, and shift right by 24 at the end.
\[
  M = \text{round}\Big(\frac{s_x s_w}{s_{\text{next}}}\cdot 2^{24}\Big),\quad
  C = \text{round}\Big(\frac{b}{s_{\text{next}}}\cdot 2^{24}\Big),\quad
  y = \big(\texttt{acc}\cdot M + C + 2^{23}\big) \gg 24.
\]
The $2^{23}$ (half of $2^{24}$) makes the shift \emph{round} instead of cutting off. Then ReLU and a
clamp to $0\dots127$ give the next layer's int8 input.

\begin{example}[: one rescale, step by step]
Take $\texttt{acc} = 2091$, a scale $s_xs_w/s_{\text{next}} = 0.0025$ and a bias that is worth $1.5$
in the next layer's units.
\begin{itemize}
  \item $M = \text{round}(0.0025\times 16{,}777{,}216) = \text{round}(41{,}943.04) = 41{,}943$.
  \item $C = 1.5\times16{,}777{,}216 = 25{,}165{,}824$.
  \item $2091\times 41{,}943 = 87{,}702{,}813$, plus $C$, plus $2^{23} = 8{,}388{,}608$, gives
        $121{,}257{,}245$.
  \item $121{,}257{,}245 \gg 24 = \lfloor 121{,}257{,}245 / 16{,}777{,}216\rfloor = \lfloor 7.23\rfloor = 7$.
\end{itemize}
Check with decimals: $2091\times0.0025 + 1.5 = 6.73$, which rounds to 7. The same answer, using only
whole numbers.
\end{example}

The $\texttt{acc}\cdot M$ product is done \emph{exactly} (with a shift-and-add loop), but only once per
neuron: 394 times per picture for the MLP, against 234{,}752 MACs. It is under 1\% of the time.

\begin{inthecode}{hw/run\_ibex.py and hw/sw/dnn.c}
\code{int\_model()} in \file{hw/run\_ibex.py} turns a trained PyTorch checkpoint into int8 weights
\code{W}, multipliers \code{M} and offsets \code{C} (with \code{SH = 24}) and writes them into a C header
\file{model.h}. \code{rescale()} and \code{relu\_q()} in \file{hw/sw/dnn.c} do the formula above on the CPU.
\end{inthecode}

\subsection{Running approximate multipliers fast on a GPU}\label{sec:factors}

A GPU is very fast at \emph{exact} matrix multiplication (many MACs at once), but it has no instruction
for ``RoBA-multiply''. Doing 10{,}000 test pictures $\times$ 400{,}000 products one by one in Python
would take far too long. The trick (in \file{roba/multipliers.py} and \file{roba/matmul.py}): every
RoBA method can be rewritten as a sum of a few \emph{exact} products of per-number pieces. For Dn1:
\[
  \text{Dn1}(x,w) = \Ar(x)\cdot w + \big(x-\Ar(x)\big)\cdot\Br(w),
\]
so a whole layer is two ordinary matrix multiplications:
$\Ar(X)\,W^{\mathsf T} + (X-\Ar(X))\,\Br(W)^{\mathsf T}$.

\begin{example}[: the piece trick on $13\times 11$]
$\Ar(13)\cdot 11 + (13-8)\cdot\Br(11) = 8\cdot 11 + 5\cdot 8 = 88+40 = 128$, the same as Dn1. A two-step
method just adds two more pieces. Only Mitchell (because of its ``if carry'' rule) cannot be split,
so it is computed element by element in chunks. That is why Mitchell's fine-tuning took about 4 times
longer than the others.
\end{example}

\subsection{Fine-tuning with the approximate multiplier}

After quantizing, we can also \textbf{fine-tune}: train a few more epochs (3 for MNIST, 5 for
Fashion-MNIST, with a small learning rate) while the network uses the approximate multiplier. The
forward pass uses the real approximate products. The backward pass (which computes how to change the
weights) pretends the multiplier is exact, because the approximate one has no usable slope. This
pretending is called the \emph{straight-through estimator}.

\begin{kid}
Fine-tuning is like practising with a slightly wrong calculator. After a while you learn how it is
wrong and adjust your own numbers so that the final answers come out right.
\end{kid}
\newpage
